\section{Quivers: the weak excluded middle is determinism}\label{sec:machines}

A \emph{quiver} $R$ is a directed graph (multiple arrows and loops allowed). A \emph{representation} of $R$ in sets gives a set $F(v)$ for each vertex and a map $F(a) : F(v)\to F(w)$ for each arrow $a : v\to w$, with no equations. A sub-representation $A$ is a family $A(v)\subseteq F(v)$ stable under the maps. Call an element $x\in F(v)$ \emph{avoiding} $A$ if no image of $x$ (along any path) lies in $A$, and \emph{attracted} by $A$ if every image of $x$ has a further image in $A$.

\begin{theorem}\label{th:quivers}
Representations of a quiver $R$ in sets form a De Morgan topos (for every representation, every sub-representation $A$ and every element $x$, $x$ is avoiding or attracted) iff every vertex of $R$ has at most one outgoing arrow. They form a Boolean topos iff $R$ has no arrows.
\end{theorem}

\begin{trace}
\item[deck] representations of $R$ $=$ presheaves on the free category $C = \mathrm{Paths}(R^{\op})$.
\item[arch] Bridge~\ref{br:dm} (De Morgan $\iff$ Ore) and Bridge~\ref{br:bool} (Boolean $\iff$ groupoid).
\item[specialise] in a free category, $f\circ h = g\circ k$ is an equality of paths, so one of $f$, $g$ is a final segment of the other. Ore on $C$ $\iff$ of any two paths ending at a vertex of $R^{\op}$, one is a final segment of the other $\iff$ no vertex of $R^{\op}$ has two incoming arrows. A free category is a groupoid iff it has no arrows.
\item[unfold] Lemma~\ref{lem:neg}: \q{$x\in\neg A$} is \q{avoiding}, \q{$x\in\neg\neg A$} is \q{attracted}.
\item[rewrite] \q{no vertex of $R^{\op}$ has two incoming arrows} $=$ \q{no vertex of $R$ has two outgoing arrows}.
\item[verify] direct proof below; Lean: %LEAN-QUIVER%
being formalised (\S\ref{sec:lean-lit}).
\item[novelty] not found in print.
\end{trace}

\direct{} If every vertex has at most one outgoing arrow, the paths from $v$ form one sequence $p_0\subset p_1\subset\cdots$, each extending the previous one. If $F(p_k)(x)\in A$, then from any image $F(p_j)(x)$ the further image $F(p_{\max(j,k)})(x)$ lies in $A$, so $x$ is attracted; otherwise it is avoiding. Conversely, let $a : v\to w$ and $b : v\to w'$ be two distinct arrows, $F(u)$ the set of paths from $v$ to $u$ (maps extend paths), $A$ the paths beginning with $a$, and $x$ the empty path at $v$. Its image along $a$ is in $A$, so it is not avoiding; its image along $b$ has no image in $A$, so it is not attracted. For the Boolean case, if $a : v\to w$ is an arrow, take $F(v) = F(w) = \{\ast\}$, $A(w) = \{\ast\}$, $A(v) = \emptyset$: the element of $F(v)$ is neither in $A$ nor avoiding.\qed

\begin{figure}[ht]
\centering
\begin{tikzpicture}[vx/.style={circle,draw,inner sep=1.5pt,minimum size=6mm,font=\small},>=Stealth]
\node[vx] (v) at (0,0) {$v$}; \node[vx,fill=red!12] (w) at (2.6,0.8) {$w$}; \node[vx] (w2) at (2.6,-0.8) {$w'$};
\draw[->] (v) -- node[above,font=\small]{$a$} (w); \draw[->] (v) -- node[below,font=\small]{$b$} (w2);
\node[align=left,font=\small] at (7.6,0) {$F(u)$ = paths from $v$ to $u$; $A$ = paths beginning with $a$;\\ the empty path at $v$ is neither avoiding nor attracted.};
\end{tikzpicture}
\caption{A branching vertex breaks the weak excluded middle.}\label{fig:branch}
\end{figure}

\inwords{} The weak excluded middle holds exactly for \emph{deterministic} quivers, those in which each vertex has a unique continuation if any: chains, cycles, and \q{in-forests} whose arrows point towards roots. One branching vertex anywhere breaks it.

\begin{corollary}\label{co:automata}
Right actions of the free monoid $\Sigma^*$ (deterministic automata over $\Sigma$, with no initial or final states) form a De Morgan topos iff $|\Sigma|\le 1$, and a Boolean topos iff $\Sigma = \emptyset$.
\end{corollary}

\begin{proof}
An action of $\Sigma^*$ is a representation of the quiver with one vertex and $|\Sigma|$ loops.
\end{proof}

\begin{trace}
\item[verify] Lean: \lean{isDeMorganTopos_singleObj_freeMonoid_unit} (one letter) and \lean{not_isDeMorganTopos_freeMonoid_two} (two letters), proved from \lean{isDeMorganTopos_singleObj_iff}.
\item[novelty] the topos of $\Sigma$-automata is standard; this property not found in print.
\end{trace}

\begin{exercise}
Which quivers among $\bullet\to\bullet\to\bullet$, $\bullet\rightleftarrows\bullet$, $\bullet\leftarrow\bullet\to\bullet$, $\bullet\to\bullet\leftarrow\bullet$ have De Morgan representation topoi? (Answer: all but the third.)
\end{exercise}
